% Requires \usepackage{graphicx} in the paper preamble.
% These are separate figure examples; select the ones appropriate for the paper.
\begin{figure}[htbp]
\centering
\includegraphics[width=\linewidth]{results/paper_figures_public/part1_01_genes_kept_at_each_filtering_step.png}
\caption{Protein-coding gene counts after annotation, expression and genotype-overlap filtering for the genome-wide cis-eQTL scan.}
\label{fig:part1-01-genes-kept-at-each-filtering-step}
\end{figure}
% In text: Figure~\ref{fig:part1-01-genes-kept-at-each-filtering-step}.

\begin{figure}[htbp]
\centering
\includegraphics[width=\linewidth]{results/paper_figures_public/part1_02_multiple_testing_results.png}
\caption{Numbers of SNP–gene pairs and genes passing nominal, genome-wide BH and Bonferroni thresholds across 6,597,508 tests.}
\label{fig:part1-02-multiple-testing-results}
\end{figure}
% In text: Figure~\ref{fig:part1-02-multiple-testing-results}.

\begin{figure}[htbp]
\centering
\includegraphics[width=\linewidth]{results/paper_figures_public/part1_03_significant_cis_eqtls_by_chromosome.png}
\caption{Genome-wide significant cis-eQTL associations summarized by autosome using BH and Bonferroni correction over the full testing family.}
\label{fig:part1-03-significant-cis-eqtls-by-chromosome}
\end{figure}
% In text: Figure~\ref{fig:part1-03-significant-cis-eqtls-by-chromosome}.

\begin{figure}[htbp]
\centering
\includegraphics[width=\linewidth]{results/paper_figures_public/part1_04_snps_tested_per_gene.png}
\caption{Distribution of tested cis-SNP counts per eligible protein-coding gene using gene bodies plus/minus 500 kb.}
\label{fig:part1-04-snps-tested-per-gene}
\end{figure}
% In text: Figure~\ref{fig:part1-04-snps-tested-per-gene}.

\begin{figure}[htbp]
\centering
\includegraphics[width=\linewidth]{results/paper_figures_public/part1_05_manhattan_strongest_cis_eqtl_per_gene.png}
\caption{Strongest cis-eQTL association per tested protein-coding gene, arranged across the 22 autosomes.}
\label{fig:part1-05-manhattan-strongest-cis-eqtl-per-gene}
\end{figure}
% In text: Figure~\ref{fig:part1-05-manhattan-strongest-cis-eqtl-per-gene}.

\begin{figure}[htbp]
\centering
\includegraphics[width=\linewidth]{results/paper_figures_public/part1_06_RPS26_association_and_expression_by_genotype.png}
\caption{Regional cis-eQTL evidence for RPS26 and expression stratified by its lead-variant genotype among matched public individuals.}
\label{fig:part1-06-RPS26-association-and-expression-by-genotype}
\end{figure}
% In text: Figure~\ref{fig:part1-06-RPS26-association-and-expression-by-genotype}.

\begin{figure}[htbp]
\centering
\includegraphics[width=\linewidth]{results/paper_figures_public/part1_07_RPS26_locuszoom.png}
\caption{Preserved RPS26 regional association visualization with LD and gene-position annotations on GRCh37.}
\label{fig:part1-07-RPS26-locuszoom}
\end{figure}
% In text: Figure~\ref{fig:part1-07-RPS26-locuszoom}.

\begin{figure}[htbp]
\centering
\includegraphics[width=\linewidth]{results/paper_figures_public/part2_pairwise_scatterplots.pdf}
\caption{Pairwise standardized disease scores among the same public reference individuals, with Pearson correlation.}
\label{fig:part2-pairwise-scatterplots}
\end{figure}
% In text: Figure~\ref{fig:part2-pairwise-scatterplots}.

\begin{figure}[htbp]
\centering
\includegraphics[width=\linewidth]{results/paper_figures_public/part2_prs_correlations.pdf}
\caption{Pearson correlations of reference-standardized disease scores in the same 489 EUR individuals.}
\label{fig:part2-prs-correlations}
\end{figure}
% In text: Figure~\ref{fig:part2-prs-correlations}.

\begin{figure}[htbp]
\centering
\includegraphics[width=\linewidth]{results/paper_figures_public/part2_prs_distributions.pdf}
\caption{Disease-specific log-odds weighted polygenic score distributions among 489 public EUR reference individuals at fixed p$\leq$0.001, clumping r$^2$=0.1 and 250 kb.}
\label{fig:part2-prs-distributions}
\end{figure}
% In text: Figure~\ref{fig:part2-prs-distributions}.

\begin{figure}[htbp]
\centering
\includegraphics[width=\linewidth]{results/paper_figures_public/part2_prs_pc_adjusted_correlations.pdf}
\caption{Pearson correlations after residualizing each standardized disease score on an intercept, PC1 and PC2 in the same people.}
\label{fig:part2-prs-pc-adjusted-correlations}
\end{figure}
% In text: Figure~\ref{fig:part2-prs-pc-adjusted-correlations}.

\begin{figure}[htbp]
\centering
\includegraphics[width=\linewidth]{results/paper_figures_public/part2_snp_overlap.pdf}
\caption{Jaccard similarities between the selected disease scoring SNP sets.}
\label{fig:part2-snp-overlap}
\end{figure}
% In text: Figure~\ref{fig:part2-snp-overlap}.

\begin{figure}[htbp]
\centering
\includegraphics[width=\linewidth]{results/paper_figures_public/part2_standardized_prs.pdf}
\caption{Disease score distributions standardized by the reference mean and sample SD.}
\label{fig:part2-standardized-prs}
\end{figure}
% In text: Figure~\ref{fig:part2-standardized-prs}.

\begin{figure}[htbp]
\centering
\includegraphics[width=\linewidth]{results/paper_figures_public/part2_reference_PCA.pdf}
\caption{PC1 and PC2 from LD-pruned autosomal EUR reference genotypes, with explained fractions from all nontrivial eigenvalues.}
\label{fig:part2-reference-PCA}
\end{figure}
% In text: Figure~\ref{fig:part2-reference-PCA}.

\begin{figure}[htbp]
\centering
\includegraphics[width=\linewidth]{results/paper_figures_public/week7_A_snp_overlap.pdf}
\caption{Percentage of unique tested cis-eQTL SNPs overlapping lifted B-cell ATAC-seq peaks, grouped by the original genome-wide significance definitions.}
\label{fig:week7-A-snp-overlap}
\end{figure}
% In text: Figure~\ref{fig:week7-A-snp-overlap}.

\begin{figure}[htbp]
\centering
\includegraphics[width=\linewidth]{results/paper_figures_public/week7_B_RPS26_accessibility.pdf}
\caption{RPS26 cis-eQTL associations with a lead-SNP marker, B-cell ATAC-seq peak track, and GENCODE v12 gene-body track on GRCh37.}
\label{fig:week7-B-RPS26-accessibility}
\end{figure}
% In text: Figure~\ref{fig:week7-B-RPS26-accessibility}.

\begin{figure}[htbp]
\centering
\includegraphics[width=\linewidth]{results/paper_figures_public/week7_C_gene_overlap.pdf}
\caption{Autosomal GENCODE v12 protein-coding genes overlapping peaks within gene bodies or gene-body plus/minus 500 kb windows.}
\label{fig:week7-C-gene-overlap}
\end{figure}
% In text: Figure~\ref{fig:week7-C-gene-overlap}.

\begin{figure}[htbp]
\centering
\includegraphics[width=\linewidth]{results/paper_figures_public/prediction_association_strength_vs_prediction.pdf}
\caption{Training-only cis-eQTL association strength versus held-out predictive R2 across the predefined pilot.}
\label{fig:prediction-association-strength-vs-prediction}
\end{figure}
% In text: Figure~\ref{fig:prediction-association-strength-vs-prediction}.

\begin{figure}[htbp]
\centering
\includegraphics[width=\linewidth]{results/paper_figures_public/prediction_atac_vs_prediction.pdf}
\caption{B-cell ATAC-seq SNP fractions and cis-window peak counts versus held-out expression predictive R2 on GRCh37.}
\label{fig:prediction-atac-vs-prediction}
\end{figure}
% In text: Figure~\ref{fig:prediction-atac-vs-prediction}.

\begin{figure}[htbp]
\centering
\includegraphics[width=\linewidth]{results/paper_figures_public/prediction_gene_characteristics_vs_prediction.pdf}
\caption{Gene length, local SNP density, training expression variability, within-gene Bonferroni-significant training SNP count and chromosome versus held-out predictive R2, with Spearman correlations.}
\label{fig:prediction-gene-characteristics-vs-prediction}
\end{figure}
% In text: Figure~\ref{fig:prediction-gene-characteristics-vs-prediction}.

\begin{figure}[htbp]
\centering
\includegraphics[width=\linewidth]{results/paper_figures_public/prediction_predictive_r2_distribution.pdf}
\caption{Distribution and evidence-stratified cumulative distributions of predictive R2 in the frozen 150-gene pilot, retaining negative values.}
\label{fig:prediction-predictive-r2-distribution}
\end{figure}
% In text: Figure~\ref{fig:prediction-predictive-r2-distribution}.

\begin{figure}[htbp]
\centering
\includegraphics[width=\linewidth]{results/paper_figures_public/prediction_single_gene_prediction.pdf}
\caption{Observed versus predicted normalized expression in 35 held-out individuals for the strongest training-only protein-coding candidate, alongside validation threshold selection.}
\label{fig:prediction-single-gene-prediction}
\end{figure}
% In text: Figure~\ref{fig:prediction-single-gene-prediction}.

\begin{figure}[htbp]
\centering
\includegraphics[width=\linewidth]{results/paper_figures_public/prediction_snp_counts_vs_prediction.pdf}
\caption{Eligible cis-SNP, training-LD clump and selected-model SNP counts versus held-out predictive R2.}
\label{fig:prediction-snp-counts-vs-prediction}
\end{figure}
% In text: Figure~\ref{fig:prediction-snp-counts-vs-prediction}.

\begin{figure}[htbp]
\centering
\includegraphics[width=\linewidth]{results/paper_figures_public/prs_exploration_disease_eqtl_overlap.pdf}
\caption{Percentage of tested disease-scoring SNPs with a BH-significant cis-eQTL, compared with harmonized and tested GWAS-background SNPs.}
\label{fig:prs-exploration-disease-eqtl-overlap}
\end{figure}
% In text: Figure~\ref{fig:prs-exploration-disease-eqtl-overlap}.

\begin{figure}[htbp]
\centering
\includegraphics[width=\linewidth]{results/paper_figures_public/prs_exploration_threshold_rank_stability.pdf}
\caption{Spearman correlations of reference PRS rankings across five predefined P+T thresholds in the same EUR individuals.}
\label{fig:prs-exploration-threshold-rank-stability}
\end{figure}
% In text: Figure~\ref{fig:prs-exploration-threshold-rank-stability}.

